% !TEX program = tectonic
% !TEX root = zwischenbericht.tex
\documentclass[type=zwischenbericht]{zukunftbau}

\setcounter{secnumdepth}{4}

\title{Zwischenbericht --- Entwerfen mit Bestand}
\subtitle{Eine offene Plattform für einen KI-unterstützten, performance-optimierten und integrativen Entwurfsprozess mit wiederverwendeten Baukomponenten}
\antragstellendeinstitution{%
  Leibniz Universität Hannover\\
  Fakultät für Architektur und Landschaft, IEK\\
  Abt.\ Nachhaltige Gebäudesysteme\\
  Projektleitung: Ueli Saluz\\
  Weitere Bearbeitung: Nikolaus Möllenhof
}
\kooperationspartner{%
  Universität der Künste Berlin\\
  Konstruktives Entwerfen und Tragwerksplanung\\
  Bearbeitung: Kinan Sarakbi
}
\aktenzeichen{10.08.18.7-25.06}
\berichtszeitraum{01.01.--10.07.2026}
\foerderzeitraum{11/2025 - 05/2027}
\footerauthors{Ueli Saluz, Kinan Sarakbi}
\makeatletter
\renewcommand{\semio@cover@zukunftbau@logo}{asset/logo/zukunft-bau.png}
\renewcommand{\semio@cover@bbsr@logo}{asset/logo/bbsr.png}
\renewcommand{\semio@cover@bundesministerium@logo}{asset/logo/gefördert-durch-bmwsb.png}
\makeatother
\kurzfassung{%
  Entwerfen mit inhomogenen Baukomponenten ist ästhetisch, historisch und ökologisch reizvoll, aber entwurfstechnisch anspruchsvoll und komplex, da funktionale, konstruktive und qualitative Abhängigkeiten direkt beim Entwerfen berücksichtigt werden müssen.
  Die im Projekt entwickelte End-To-End-Plattform, welche gleichzeitig Katalog-, Entwurfs- und Performancebeurteilungs-Tool beinhaltet, ermöglicht es, in dieser Situation effektiv Lösungen zu entwickeln.
  Das Performance-Assessment hinsichtlich Energie und Tragwerk erlaubt frühe multidisziplinäre Optimierung.
  Mithilfe einer KI-Assistenz können auch große Bestandspools, in welchen es explosionsartig viele Kombinationsmöglichkeiten gibt, im Entwurfsprozess berücksichtigt werden.
  Die Assistenz schlägt für die jeweilige Entwurfssituation passende Baukomponenten und Verknüpfungen vor.
  Basierend auf einer semantischen Repräsentation der Baukomponenten und des komponentenbasierten Entwurfes wird die generative Stärke der vortrainierten Large-Language-Models direkt für den Entwurf genutzt.
  Über eine Framework-unabhängige Schnittstellendefinition wird es bestehenden Bauteilbörsen ermöglicht, ihre Baukomponenten nahtlos maschinenlesbar zu veröffentlichen, sodass die Integrierbarkeit in die Plattform gewährleistet ist.
  Das Schema dieser Schnittstellendefinition beinhaltet neben Fotos und den üblichen Metadaten wie Ort und Preis auch 3D-Modelle und Performance-Kennwerte wie das Treibhauspotential, welche direkt von den Tools verwendet werden.
  Ein Modellierungstool, welches moderne Bild-zu-3D-KI-Modelle benutzt, reduziert den Aufbereitungsaufwand der Bauteilbörsen.
  Durch die open-source Veröffentlichung der gesamten Plattform, deren Dokumentation und den Aufbau einer Community werden die Langlebigkeit und das Vertrauen gestärkt.
  Das Self-Hosting ermöglicht einerseits eine Breitenwirkung, wodurch der zukünftige Unterhalt und neue Featureentwicklung effizient geteilt wird.
  Andererseits bleibt die Datensouveränität gewährleistet.%
}

\makeatletter
\renewcommand{\semio@glossary@header@definition}{Verwendung im Bericht}
\makeatother

\GlossaryDefine{Entwurfsfähigkeit}{Fähigkeit eines wiederverwendeten Bauteils, in frühen Entwurfsphasen gesucht, verglichen, bewertet und einer möglichen Rolle in einer Variante zugeordnet zu werden.}
\GlossaryDefine{Bauteil}{Konkretes physisches Element aus einem Bestand, z.\,B. Stütze, Träger, Platte oder Zuschnittfragment.}
\GlossaryDefine{Baukomponente}{Übergeordneter Begriff für einzelne Bauteile, Bauteilgruppen oder Pakete, sofern sie für den Entwurf relevant sind.}
\GlossaryDefine{Bestandsquelle}{Quelle wiederverwendbarer Baukomponenten, z.\,B. Bauteilbörse, Rückbauinventar, Materialkataster, Projektkatalog oder direkte Bestandsaufnahme.}
\GlossaryDefine{Bauteilbörse}{Plattform oder Organisation, die wiederverwendbare Baukomponenten sichtbar macht, vermittelt, verkauft oder dokumentiert.}
\GlossaryDefine{Ebene A / Ebene B}{Ebene A bezeichnet beschaffungsorientierte Angebotsdaten, etwa Standort, Preis oder Verfügbarkeit. Ebene B bezeichnet qualifizierte, planungsrelevante Bauteildaten für Entwurf, Vorbewertung und Nachweisführung.}
\GlossaryDefine{Planungsrelevanter Bauteildatensatz}{Strukturierter Datensatz mit Geometrie, Material, Zustand, Herkunft, Verfügbarkeit, möglicher Entwurfsrolle, Bewertungsannahmen und Informationszustand.}
\GlossaryDefine{Qualifizierung}{Überführung unvollständiger oder beschaffungsorientierter Bestandsdaten in planungsrelevante Informationen.}
\GlossaryDefine{Datenreifegrad}{Einschätzung, wie belastbar einzelne Informationen zu einer Baukomponente sind.}
\GlossaryDefine{Informations- und Nachweiszustände}{Vier Zustände zur Kennzeichnung einzelner Angaben: vorhanden, abgeleitet, angenommen und offen.}
\GlossaryDefine{Entwurfsrolle}{Mögliche Funktion eines Bauteils im neuen Entwurf, z.\,B. als Träger, Stütze, Platte, Wandscheibe oder Zuschnittfragment.}
\GlossaryDefine{Typologiefamilie}{Ordnungskategorie für Baukomponenten nach Rolle und Informationsbedarf, z.\,B. Linien-, Flächen-, Rahmen-, Stütz- oder Zuschnittfragmente.}
\GlossaryDefine{Digitales Bauteilobjekt}{Digitale, entwurfsfähige Repräsentation einer realen Baukomponente mit Geometrie, Parametern, Rollen, Anschlussinformationen und Bewertungsannahmen.}
\GlossaryDefine{Generator}{Parametrischer Mechanismus, der aus qualifizierten Eingaben ein digitales Bauteilobjekt erzeugt.}
\GlossaryDefine{Bauteilportal}{Eingangs- und Qualifizierungsebene des Systems. Hier werden Baukomponenten erfasst, importiert, ergänzt und strukturiert.}
\GlossaryDefine{Bauteilkatalog}{Such-, Sicht- und Auswahlebene des Systems. Hier werden qualifizierte Baukomponenten auffindbar, filterbar und vergleichbar gemacht.}
\GlossaryDefine{Playground}{Entwurfs- und Bewertungsebene des Systems. Hier werden Baukomponenten zu Varianten kombiniert und ökologisch sowie tragwerklich vorbewertet.}
\GlossaryDefine{Aggregator}{Technische Bezeichnung für den Playground innerhalb des semio-basierten Entwurfswerkzeugs.}
\GlossaryDefine{Wizard}{Geführte Eingabeoberfläche zur schrittweisen Erfassung und Qualifizierung von Bauteildaten.}
\GlossaryDefine{Test-Case}{Realer oder realitätsnaher Anwendungsfall, an dem die entwickelte digitale Kette geprüft wird.}

\begin{document}

\makecoverpages

\maketableofcontents

\makeworkpackages{%
  \SemioTableThree{%
    \SemioTableHeaderRow{AP & Partner & Schwerpunkt im Berichtszeitraum}
    \SemioTableRow{AP-Erfahrung & Leibniz Universität Hannover & Recherche, Interviews, User Stories}
    \SemioTableRow{AP-Plattform & Leibniz Universität Hannover & Generatoren, Bauteileinspeisung}
    \SemioTableRow{AP-Tool & Universität der Künste Berlin & HID, Lebenszyklusanalyse, Tragwerksanalyse, KI-Assistenz}
    \SemioTableRow{AP-Validierung & Leibniz Universität Hannover & Test-Case, Entwurfsgrammatik, Entwicklung}
  }%
}

\section{Ergebnisse}

\subsection{Forschungsfragen und Vorgehen}

Das Vorhaben untersucht, wie wiederverwendete Baukomponenten in frühen Entwurfsphasen skalierbar planungsfähig gemacht werden können. Ausgangspunkt ist die Beobachtung, dass die Zahl der Bauteilbörsen, Rückbauinventare oder Materialbanken (siehe Anlagen) stärker wächst als die Zahl der realisierten Projekte (siehe Anlage Projekte). Die Hürden, welchen es zu überwältigen gilt, um mit wiederverwendeten Baukomponenten zu bauen, ist lang (siehe Anlage Hürden und Handbuch Abbau Aufbau). Unser Beitrag ist es, die effiziente, integrative und KI-gestützte Vorplanung für Entwürfe zu ermöglichen.

Vier Forschungsfragen leiten die Arbeit:

\begin{enumerate}
\item Wie lassen sich heterogene Baukomponenten verallgemeinern/systematisieren, damit sie möglichst flexibel in neuen Entwurfssituationen eingesetzt werden können?
\item Wie können Bauteilbörsen mit minimalem Aufwand planungsfähige Daten bereitstellen?
\item Wie können gleichzeitig Entwurfs-, LCA- und Tragwerksmodelle erstellt/modifiziert/evaluiert werden?
\item Wie kann KI, insbesondere LLMs, während dem Entwurfsprozess unterstützen?
\end{enumerate}

Der Zwischenbericht folgt dabei der Struktur der Arbeitspakete. AP-Erfahrung leitet aus der Recherche und den Interviews User Stories ab. Die User Stories definieren die fachlichen Anforderungen für die Plattform und für das Tool. AP-Plattform und AP-Tool übersetzen diese Anforderungen in digitale Systemkomponenten. AP-Validierung prüft die entwickelten Ansätze anhand eines Test-Cases.

\subsection{Arbeitspakete}

\subsubsection{AP-Erfahrung}

\paragraph{Recherche}
Die Recherche ergab, dass wiederverwendete Bauteile in bestehenden Quellen zwar auffindbar, aber selten unmittelbar entwurfsfähig sind. Entwurfsfähig meint hier, dass ein Bauteil in frühen Entwurfsentscheidungen gesucht, verglichen, bewertet und einer möglichen Rolle in einer Variante zugeordnet werden kann.

Die zentrale Lücke liegt zwischen Sichtbarkeit und Planbarkeit. Bauteilbörsen, Rückbauinventare und Materialbanken liefern meist Angaben zu Verfügbarkeit, Standort, Zustand oder Maßen. Für frühe Entwurfsentscheidungen fehlen häufig belastbare Angaben zu möglicher Rolle, Nachweisstand, Anschlussfähigkeit, Bewertungsannahmen und Unsicherheiten.

Daraus wurde die Unterscheidung zwischen beschaffungsorientierten Daten (Ebene~A) und planungsrelevanten Bauteildatensätzen (Ebene~B) abgeleitet. Ein Bauteil wird planungsfähig, wenn es qualifiziert, typologisch eingeordnet und mit einem vorläufigen Bewertungs- und Nachweisstand versehen ist.

\begin{SemioNest}
\paragraph{Bauteilbörsen}
Die Auswertung von 51 öffentlich zugänglichen Bauteilbörsen und Angebotskanälen zeigt vier wiederkehrende Angebotslogiken: Depot-Shop, Marktplatz, Harvest-Katalog und Vermittlungsplattform (vgl.\ Anlage Bauteilbörsen). Die Auswertung folgt der Außensicht auf die Plattformen und erfasst, was Angebotsseiten, Kataloge und Produktkarten öffentlich preisgeben. Die Typen unterscheiden sich vor allem durch die Qualität, Stabilität und Anschlussfähigkeit ihrer Daten. Für das Vorhaben ist daher nicht die Plattformart selbst entscheidend, sondern ob aus einem Angebot ein planungsfähiger Bauteildatensatz entstehen kann.

\begin{SemioNest}
\paragraph{Depot-Shop}
weist tendenziell die stabilste Datenlage auf, weil Bauteile häufig bereits ausgebaut, gelagert oder gesichtet sind. Relevante Angaben sind Lagerort, Reservierbarkeit, Maße, Zustand und Verfügbarkeit. Die Angaben reichen meist aus für eine 3D-Modellierung, allerdings ist eine Vorbewertung und Nachweisführung nur aus den vorhandenen Daten nicht möglich.

\paragraph{Marktplatz}
weist potentiell das höchste Bauteilangebot auf, führt aber häufig zu uneinheitlicher Datenqualität. Die Detailtiefe fluktuiert am meisten, da sie stark von den jeweiligen Anbieter:innen abhängt. Da meist keine Richtlinien zur Datenerfassung vorliegen oder eingehalten werden, ergibt sich ein Bedarf an Datenprüfung, Ergänzung und Normalisierung.

\paragraph{Harvest-Katalog}
weist das größte Optimierungspotential auf, da gegebenfalls bedarfsgerecht rückgebaut werden kann, allerdings auf Kosten der Skalierbarkeit, erhöhter Unsicherheit und damit risikobehafteter. Zusatzinformation ist das Zeitfenster, Kosten und Aufwand für den Rückbau. Terminplanung, Logistik und Alternativen müssen insbesondere berücksichtigt werden.

\paragraph{Vermittlungsplattform}
stellt Kontakte zwischen Besitzenden, Suchenden und Fachakteur:innen her. Sie kann Wiederverwendungsprozesse anstoßen, führt aber häufig keine vollständigen Bauteildatensätze. Für die Planung wird sie erst anschlussfähig, wenn aus der Vermittlung ein strukturierter Datensatz entsteht.
\end{SemioNest}

Die Auswertung zeigt, dass die meisten bestehenden Angebotsquellen Bauteile sichtbar machen, aber selten vollständige Entscheidungsgrundlagen für Entwurf, Bewertung und Nachweisführung liefern; dieser Befund deckt sich mit der Literatur zu unstrukturierten Märkten und schwacher Rückverfolgbarkeit tragender Bauteile (Tingley et al.\ 2017; Sivers et al.\ 2025). Benötigt werden Angaben zu Geometrie, Zustand, Herkunft, Verfügbarkeit, möglicher Rolle, Nachweisstand, Datenquelle und Unsicherheit. Inventarstandards wie Swiss-Inv bestätigen diese Richtung, weil sie auf harmonisierte Bauteildaten und bessere Übertragbarkeit zwischen Plattformen, Planung und Dokumentation zielen (Cirkla o.\,J.). Für das Vorhaben folgt daraus: Die Eingabe- und Qualifizierungsebene muss Angebotsdaten in qualifizierte Bauteildatensätze überführen. Dazu gehören ein Datenreifegrad, ein Verfügbarkeitsstatus, ein Nachweisstatus und die Kennzeichnung fehlender oder unsicherer Angaben.

\paragraph{Wiederverwendungsprozesse}
Die Recherche zu Wiederverwendungsprozessen zeigt, dass Entwurfsfähigkeit früh im Prozess entsteht oder verloren geht. Entscheidend ist, ob Informationen aus Rückbau, Inventarisierung und Bauteilhandel rechtzeitig in die Planung gelangen. Digitale ReUse-Workflows wie der D5-Workflow (De Wolf et al.\ 2024) oder der buildingSMART-Use-Case „Decommissioning to ReUse" bestätigen, dass Wiederverwendung auf durchgängigen Datenflüssen zwischen Erfassung, Rückbau, Verteilung, Planung und Wiedereinbau beruht. Für das Vorhaben ist daraus vor allem relevant, dass ein Bauteil bereits in der Vorplanung mit einem nachvollziehbaren Informationsstand vorliegen muss. Planungsrelevant wird es erst, wenn Geometrie, Zustand, Verfügbarkeit, mögliche Rolle und offene Prüfungen erkennbar sind.

\begin{SemioNest}
\paragraph{Gestaltung}
Die Recherche bestätigt, dass Wiederverwendung die Richtung des Entwerfens verändert. Beim Entwerfen mit wiederverwendeten Bauteilen wirken sich vorhandene Längen, Querschnitte, Zustände, Anschlüsse und Verfügbarkeiten auf die frühe Formfindung aus; dieser Zusammenhang ist als „form follows availability" beschrieben (Brütting et al.\ 2019). Darauf aufbauend lassen sich zwei Strategien unterscheiden: das Entwerfen aus dem Bestand, das ausschließlich vorhandene Bauteile verwendet, und das Entwerfen mit dem Bestand, das wiederverwendete Bauteile einsetzt, wo es möglich ist, und mit Neuware ergänzt, wo es nötig ist (Bertin et al.\ 2020). Daraus folgt, dass Bauteile über ihre Produktbezeichnung hinaus beschrieben werden müssen. Ein Bauteil ist entwurfsrelevant, wenn es in einer Variante eine nachvollziehbare Rolle übernehmen kann: als Träger, Stütze, Platte, Wandscheibe, Rahmenteil oder Zuschnittfragment. Die Suche und Auswahl muss daher mögliche Entwurfsrollen, geometrische Einschränkungen und Bewertungszustände berücksichtigen.

\paragraph{Nachweise}
Die Recherche zeigt, dass Nachweise in frühen Entwurfsphasen nicht nur als abgeschlossen oder fehlend verstanden werden können. Wiederverwendete Bauteile liegen häufig mit unterschiedlichen Informationsständen vor. Einige Angaben sind belegt, andere aus Bestandsunterlagen oder Erfassung abgeleitet, weitere beruhen auf Annahmen oder bleiben offen. Der Datenreifegrad beschreibt, wie belastbar einzelne Angaben sind. Er wird im Vorhaben über vier Nachweiszustände gefasst: vorhanden, abgeleitet, angenommen und offen. Diese Zustände betreffen geometrische, materielle, tragwerkliche und ökologische Informationen. Geometrische Nachweise umfassen Maße, Toleranzen, Schnittflächen, Verformungen und Anschlussstellen. Materielle und tragwerkliche Nachweise betreffen Materialkennwerte, Zustand, Vorschädigungen, Lastabtrag, Verbindungsmöglichkeiten und erforderliche Prüfungen. Für die ökologische Vorbewertung sind Systemgrenze, Wiederverwendungspfad, Transport, Instandsetzung, Allokationsregel und substituiertes Primärmaterial relevant; diese Angaben liegen in der Vorplanung häufig nur als Annahmen vor und müssen sichtbar bleiben. Für den Umweltnachweis gilt dabei das Prinzip der projektspezifischen Substitution: Der Vorteil eines wiederverwendeten Bauteils bemisst sich am ersetzten Primärmaterial im konkreten Projekt (Abschnitt~1.2.3.3). Für das Werkzeug folgt daraus eine klare Funktion. Es soll nicht die abschließende Fachplanung ersetzen. Es soll sichtbar machen, auf welcher Datengrundlage eine Variante bewertet wird und welche Prüfungen im weiteren Prozess erforderlich bleiben.
\end{SemioNest}

\paragraph{Bauteiltypologien}
Aus der Recherche wurde eine typologische Arbeitsordnung abgeleitet. Sie ordnet Bauteile nach ihrer möglichen Rolle im neuen Entwurf und verbindet Erfassung, Suche, digitale Objektbildung und Vorbewertung. Die Typologie beschreibt nicht nur Bauteilarten. Sie legt fest, welche Informationen je Bauteilfamilie benötigt werden. Polymorphie bezeichnet hier keine freie Umdeutung, sondern eine prüfbare Rollenzuordnung. Ein Bauteil kann eine andere Rolle übernehmen als im ursprünglichen Gebäude. Diese Möglichkeit hängt von Orientierung, Geometrie, Material, Zustand, Anschluss, Schnittführung und Nachweisstand ab.

\begin{SemioNest}
\paragraph{Linienelemente}
stabförmige Bauteile mit dominanter Längsrichtung. Für die Erfassung sind Länge, Querschnitt, Material, Zustand, Anschlussenden und mögliche Tragrichtung relevant. Beispiele sind Träger, Balken, Unterzüge und Riegel. Für den Entwurf ist entscheidend, ob Geometrie, Tragrolle und Anschlussbedingungen hinreichend beschreibbar sind; Einschränkungen entstehen durch vorhandene Längen, Querschnitte, Anschlussstellen und mögliche Vorschädigungen.

\paragraph{Flächenelemente}
platten- oder scheibenförmige Bauteile. Relevante Angaben sind Länge, Breite, Dicke, Spannrichtung, Bewehrungs- oder Schichtaufbau, Öffnungen, Schnittkanten und Randbedingungen. Für die Wiederverwendung ist entscheidend, ob ein Flächenelement horizontal, vertikal oder aussteifend eingesetzt werden kann; die mögliche Rolle hängt von Geometrie, Lagerung, Schnittkanten, Öffnungen, Materialaufbau und Nachweisstand ab.

\paragraph{Rahmenelemente}
gefügte Teilstrukturen, deren Verhalten sich aus Verbindung, Geometrie und Steifigkeit des Gesamtsystems ergibt. Relevante Angaben sind Verbindungstyp, Zerlegbarkeit, Steifigkeit, Transportierbarkeit und mögliche Weiterverwendung als Einheit oder als Einzelteile. Hier entsteht eine frühe Entscheidungsfrage: Der Rahmen kann als zusammenhängendes Element erhalten bleiben oder in Einzelteile zerlegt werden. Diese Entscheidung beeinflusst Erfassung, Transport, Nachweisführung und spätere digitale Objektbildung.

\paragraph{Stützelemente}
primär druckbeanspruchte Bauteile. Sie werden als eigene Familie geführt, weil Knicken, Lagerung, Imperfektionen und Anschlussbedingungen andere Nachweisfragen erzeugen als biegebeanspruchte Linienelemente. Für die Erfassung sind Querschnitt, Länge, Material, Zustand, Anschlussflächen, mögliche Exzentrizitäten und Lagerungsbedingungen relevant. Eine Umnutzung als liegendes Linienelement ist möglich, erfordert aber eine neue Rollenzuordnung und einen entsprechenden Nachweisstand.

\paragraph{Zuschnittfragmente}
aus größeren Monolithteilen herausgelöste Teilstücke. Ihre spätere Rolle entsteht durch Schnittführung, Restgeometrie, Bewehrungsverlauf, Kantenqualität und Anschlussmöglichkeit; diese Familie ist für die im Vorgängervorhaben (siehe Abbau Aufbau) untersuchten Verbund-Monolithteile zentral. Zuschnittfragmente liegen nicht als vollständige Produkte vor, sondern entstehen durch eine Entscheidung im Rückbau- oder Aufbereitungsprozess. Deshalb müssen Schnittlinien, verbleibende Bewehrung, mögliche Beschädigungen, Transportierbarkeit und neue Anschlussflächen dokumentiert werden. Für das Vorhaben ist diese Familie besonders relevant, weil sie zeigt, dass ein wiederverwendetes Bauteil nicht immer als fertiges Produkt vorliegt; es kann erst durch Schnittführung, Orientierung und neue Rolle als digitales Bauteilobjekt entstehen.
\end{SemioNest}

Die Typologie ist damit ein Ergebnis der Recherche und zugleich ein Arbeitsmittel für die Systementwicklung. Sie legt fest, welche Informationen bei der Erfassung abgefragt werden, welche Filter in der Such- und Vergleichsebene sinnvoll sind, welche digitale Objektbildung benötigt wird und welche Vorprüfungen an eine gewählte Rolle gebunden sind. Die Recherche zeigt insgesamt, dass wiederverwendete Bauteile an mehreren Stellen sichtbar werden: in Bauteilbörsen, Rückbauinventaren, Projektkatalogen und Materialbanken. Für frühe Entwurfsentscheidungen reicht diese Sichtbarkeit nicht aus. Planungsfähig werden die Bauteile erst durch qualifizierte Daten, typologische Einordnung, Nachweiszustände und vorläufige Bewertbarkeit.
\end{SemioNest}

\paragraph{Interviews}
\begin{SemioNest}
\paragraph{18 Interviews}
Als breite empirische Grundlage nutzt das Vorhaben die Studie „ReUse in der Bauindustrie stärken" (Cirkla/BaselCircular 2025): 18 qualitative Fachinterviews, in 13 Hemmnis-Gruppen geordnet und in einer Umfrage mit 94 Fachleuten nach Wirksamkeit bewertet. Sie erfüllt zugleich die für Anlage~8 geforderte Stellungnahme zu Ergebnissen Dritter und eine Auflage des Zuwendungsgebers, der eine breitere empirische Basis verlangt hatte; die eigenen Fachgespräche erweitern sie gezielt (Abschnitt~1.2.1.2.2).

Die stärksten Hemmnisse liegen laut Studie nicht im Material, sondern im Prozess: Kosten und fehlende Akzeptanz führen die Rangfolge an, Informationsdefizite sowie Haftungs- und Versicherungsfragen bilden ein mittleres Niveau. Für den Entwurf mit wiederverwendeten Bauteilen sind gerade diese mittleren Gruppen ausschlaggebend, denn sie entscheiden, ob ein Bauteil planbar wird: fehlende Angaben zu Verbindungen, Lebensdauer und Tragfähigkeit zum Entwurfszeitpunkt, fehlende anerkannte Prüf- und Nachweiswege und die zu späte Berücksichtigung von ReUse in der Planung.

An diesen Punkten setzen die technologischen Ansätze an, die die Studie als wirksam bewertet: die frühe Integration von ReUse in die Planung und deren digitale Unterstützung, etwa die Kopplung von Inventarisierungsstandards mit Ausschreibungssystemen und ein ReUse-Faktor in frühen Ökobilanzen. Gegen fehlende Bauteilinformationen wirken Materialpässe wie Madaster, gegen die mangelnde Interoperabilität zwischen Plattformen, Ausschreibungen und Planungswerkzeugen der Datenstandard Swiss-Inv. Die Studie hält zugleich fest, dass diese Werkzeuge noch nicht flächendeckend etabliert sind. Für das Vorhaben bestätigt dies die Annahme, dass digitale Unterstützung dort ansetzen muss, wo Bauteilinformationen, Planungsentscheidungen und Bewertungsverfahren bisher nicht durchgängig verbunden sind.

\paragraph{Fehlende Sichtweisen}
Die 18 Interviews der Stärken-Studie decken die Breite der Baupraxis ab, lassen aber Fachperspektiven unterbelichtet, die für das Entwerfen mit Bestand zentral sind. Diese Lücke schließt das Vorhaben mit eigenen Experteninterviews, die über die Laufzeit geführt werden. Im Berichtszeitraum liegen zwei ausgewertete Gespräche vor; weitere sind angesetzt. Die eigenen Gespräche wiederholen die Sekundärstudie nicht, sondern vertiefen gezielt Perspektiven, die für Plattform, Tool und Validierung besonders relevant sind.

\begin{SemioNest}
\paragraph{Digitale Gestaltung}
Ein Gespräch mit einem Experten für computergestützten Entwurf richtete den Blick darauf, wie sich unregelmäßige, vorgeprägte Bauteile digital erfassen und im Entwurf verarbeiten lassen (vgl.\ Anlage Interviewauswertung, digitale Gestaltung). Statt Material nur zu verwalten, gehe es darum, es zu „lesen" und in ein digitales Materialobjekt zu überführen, dessen Unregelmäßigkeit als gestalterische Qualität erhalten bleibt. Für ein Werkzeug folgt daraus, Entscheidungen unter Unsicherheit zu ermöglichen und Bauteile ihrer möglichen Rolle im Entwurf zuzuordnen, statt sie nur zu katalogisieren.

\paragraph{Kreislaufwirtschaft}
Raoul Bunschoten hob die fehlende Kontinuität entlang der Wertschöpfungskette hervor: Materialherkunft, Planung, Rückbau und Wiederverwendung sind heute unzureichend verbunden (vgl.\ Anlage Interviewauswertung Bunschoten). Ein digitales Werkzeug solle zudem keine fertigen Lösungen erzeugen, sondern Entscheidungsräume öffnen und Entwurfsvarianten mit ihren zugrunde liegenden Annahmen und Unsicherheiten vergleichbar machen. Es wirke daher weniger als abgeschlossene Anwendung denn als verbindende Struktur zwischen Datenquellen, Bewertungsverfahren und Beteiligten.

\paragraph{Tragwerksplanung}
\begin{Blockquote}
\textit{[Platzhalter --- Auswertung des Experteninterviews zur Tragwerksplanung folgt; das Gespräch ist angesetzt. Erwarteter Fokus: strukturelle Nachweisführung und Prüfpfade bei wiederverwendeten tragenden Bauteilen. Geschätzter Umfang ca.\ 600 Zeichen.]}
\end{Blockquote}
\end{SemioNest}
\end{SemioNest}

\paragraph{User Stories}
Aus den Befunden von Recherche und Interviews wurden strukturierte Anforderungen abgeleitet (vgl.\ Anlage User-Stories und Anforderungsmapping). Die User Stories übersetzen die zuvor beschriebenen Befunde in funktionale Anforderungen und bilden damit die Brücke zwischen AP-Erfahrung und den nachfolgenden Arbeitspaketen AP-Plattform, AP-Tool und AP-Validierung. In einem erfahrungsbasierten Prozess entstanden 36 Anforderungskarten aus vier Akteursperspektiven, verdichtet zu 35 konsolidierten User Stories im Format „Als [Akteur:in] möchte ich [Funktion], damit [Nutzen]". Die vier Perspektiven greifen mit unterschiedlichem Bedarf auf dasselbe Bauteil zu:

\begin{SemioNest}
\paragraph{Mitarbeiter:in einer Bauteilbörse}
(10) --- Bauteildaten erfassen, gruppieren und mit Status, Herkunft, Nachweisen und Unsicherheiten versehen.

\paragraph{Architekt:in}
(12) --- Bauteile suchen, gestalterisch und materiell bewerten, in den Entwurf überführen.

\paragraph{Tragwerksplaner:in}
(6) --- strukturelle Eigenschaften filtern, Informationslücken erkennen, Nachweise definieren, Varianten vergleichen.

\paragraph{Energieberater:in}
(7) --- ökologische und energetische Kennwerte je Bauteil bereitstellen und in die Variantenbewertung einbringen.
\end{SemioNest}

Das Anforderungsmapping übersetzt die Stories in fünfzehn funktionale Cluster (REQ-01 bis REQ-15) und ordnet sie drei Systembereichen zu, die einen durchgängigen Arbeitsprozess bilden. Das \textbf{Bauteilportal} erfasst, qualifiziert und prüft Bauteildaten samt Nachweisen und Unsicherheiten; der \textbf{Bauteilkatalog} macht die qualifizierten Bauteile einheitlich such-, vergleich- und bewertbar; der \textbf{Aggregator} kombiniert ausgewählte Bauteile zu Varianten und bewertet diese tragwerklich und ökologisch. Die Schnittstellen zählen nicht doppelt, sondern gelten als Übergaben vom Portal in den Katalog und vom Katalog in den Aggregator.

\subsubsection{AP-Plattform}

\paragraph{Generatoren}

\paragraph{Bauteileinspeisung}
\begin{SemioNest}
\paragraph{UI}
\begin{SemioNest}
\paragraph{Wizard}
\end{SemioNest}
\paragraph{API}
\begin{SemioNest}
\paragraph{REST}
\end{SemioNest}
\end{SemioNest}

\subsubsection{AP-Tool}

\paragraph{Human Interface Design}

\paragraph{Lebenszyklusanalyse}

\paragraph{Tragwerksanalyse}

\paragraph{KI Assistenz}

\subsubsection{AP-Validierung}

\paragraph{Test-Case}
\begin{SemioNest}
\paragraph{Entwurfsgrammatik}
\paragraph{Entwicklung}
\end{SemioNest}

\section{Projektstand}

\subsection{Arbeitspakete}

\subsubsection{AP-Erfahrung}

\paragraph{Neuer Sekundärbefund}
\begin{SemioNest}
\paragraph{Geplante Interviews}
\begin{SemioNest}
\paragraph{I-Usability}
\begin{SemioNest}
\paragraph{I-U-Plattform}
\paragraph{I-U-Tool}
\end{SemioNest}
\paragraph{I-Validierung}
\begin{SemioNest}
\paragraph{I-V-Bauteilbörse}
\paragraph{I-V-Architekt:in}
\paragraph{I-V-Tragwerksplaner:in}
\paragraph{I-V-Energieberater:in}
\end{SemioNest}
\end{SemioNest}
\end{SemioNest}

\subsubsection{AP-Plattform}

\paragraph{Eingabe}

\paragraph{Schnittstelle}

\subsubsection{AP-Tool}

\paragraph{Human Interface Design}

\paragraph{Lebenszyklusanalyse}

\paragraph{Tragwerksanalyse}

\paragraph{KI Assistenz}

\subsubsection{AP-Validierung}

\paragraph{Test-Case}
\begin{SemioNest}
\paragraph{Entwurfsgrammatik}
\paragraph{Entwicklung}
\end{SemioNest}

\paragraph{Workshop}
\begin{SemioNest}
\paragraph{Vorbereitung}
\paragraph{Durchführung und Auswertung}
\end{SemioNest}

\section{Mittelverwendung}

\subsection{Personalkosten}

\subsubsection{Leibniz Universität Hannover}
\begin{SemioNest}
\paragraph{Ueli Saluz}
\paragraph{Nikolaus Möllenhof}
\end{SemioNest}

\subsubsection{Universität der Künste Berlin}
\begin{SemioNest}
\paragraph{Kinan Sarakbi}
\end{SemioNest}

\subsection{Material}

\subsection{Leistungen Dritter}

\subsubsection{Softwarelizenzen}

\subsubsection{Dienstleistungen}
\begin{SemioNest}
\paragraph{Cloud - Infrastructure-as-a-Service}
\paragraph{Nutzung vortrainierter KI-Modelle}
\paragraph{Externe Softwareentwicklung}
\begin{SemioNest}
\paragraph{Server}
\paragraph{Fotogrammetrie}
\paragraph{User-Interface}
\paragraph{Anbindungen Bauteilbörsen}
\end{SemioNest}
\end{SemioNest}

\subsection{Reisekosten}

\section{Ergebnisverwertung}

\subsection{Repositorien}

\subsubsection{GitHub}

\subsection{Websiten}

\subsubsection{Tool}

\subsubsection{Dokumentation}

\subsection{Community}

\subsubsection{Discourse}

\subsection{Forschung}

\subsubsection{Zenodo}

\subsubsection{Journalpaper}

\appendix
\section{Verzeichnisse}
\listofglossaries

\end{document}
